\begin{lemma}
\label{lemma-cohomological-dimension-proper}
Let $f : X \to Y$ be a proper morphism of schemes
all of whose fibres have dimension $\leq n$.
Then for any abelian torsion sheaf $\mathcal{F}$ on $X_\etale$
we have $R^qf_*\mathcal{F} = 0$ for $q > 2n$.
\end{lemma}

\begin{proof}
We will prove this by induction on $n$ for all proper morphisms.

\medskip\noindent
If $n = 0$, then $f$ is a finite morphism
(More on Morphisms, Lemma \ref{more-morphisms-lemma-characterize-finite})
and the
result is true by Proposition \ref{proposition-finite-higher-direct-image-zero}.

\medskip\noindent
If $n > 0$, then using Lemma \ref{lemma-proper-base-change-stalk}
we see that it suffices to prove $H^i_\etale(X, \mathcal{F}) = 0$
for $i > 2n$ and $X$ a proper scheme, $\dim(X) \leq n$
over an algebraically closed field $k$
and $\mathcal{F}$ is a torsion abelian sheaf on $X$.

\medskip\noindent
If $n = 1$ this follows from Theorem \ref{theorem-vanishing-curves}.
Assume $n > 1$. By Proposition \ref{proposition-topological-invariance}
we may replace $X$ by its reduction.
Let $\nu : X^\nu \to X$ be the normalization.
This is a surjective birational finite morphism
(see Varieties, Lemma \ref{varieties-lemma-normalization-locally-algebraic})
and hence an isomorphism over a dense open $U \subset X$
(Morphisms, Lemma \ref{morphisms-lemma-birational-birational}).
Then we see that
$c : \mathcal{F} \to \nu_*\nu^{-1}\mathcal{F}$ is injective
(as $\nu$ is surjective) and an isomorphism over $U$.
Denote $i : Z \to X$ the inclusion of the complement of $U$.
Since $U$ is dense in $X$ we have $\dim(Z) < \dim(X) = n$.
By Proposition \ref{proposition-closed-immersion-pushforward} have
$\Coker(c) = i_*\mathcal{G}$ for some
abelian torsion sheaf $\mathcal{G}$ on $Z_\etale$.
Then $H^q_\etale(X, \Coker(c)) = H^q_\etale(Z, \mathcal{G})$
(by Proposition \ref{proposition-finite-higher-direct-image-zero}
and the Leray spectral sequence) and by induction hypothesis we conclude that
the cokernel of $c$ has cohomology in degrees $\leq 2(n - 1)$.
Thus it suffices to prove the result for $\nu_*\nu^{-1}\mathcal{F}$.
As $\nu$ is finite this reduces us to showing
that $H^i_\etale(X^\nu, \nu^{-1}\mathcal{F})$ is
zero for $i > 2n$. This case is treated in the next paragraph.

\medskip\noindent
Assume $X$ is integral normal proper scheme over $k$ of dimension $n$.
Choose a nonconstant rational function $f$ on $X$. The graph
$X' \subset X \times \mathbf{P}^1_k$ of $f$ sits into a diagram
$$
X \xleftarrow{b} X' \xrightarrow{f} \mathbf{P}^1_k
$$
Observe that $b$ is an isomorphism over an open subscheme
$U \subset X$ whose complement is a closed subscheme
$Z \subset X$ of codimension $\geq 2$. Namely, $U$ is the
domain of definition of $f$ which contains all codimension $1$
points of $X$, see
Morphisms, Lemmas \ref{morphisms-lemma-rational-map-from-reduced-to-separated}
and \ref{morphisms-lemma-extend-across}
(combined with Serre's criterion for normality, see
Properties, Lemma \ref{properties-lemma-criterion-normal}).
Moreover the fibres of $b$ have dimension $\leq 1$ (as closed subschemes
of $\mathbf{P}^1$). Hence $R^ib_*b^{-1}\mathcal{F}$ is nonzero only
if $i \in \{0, 1, 2\}$ by induction. Choose a distinguished triangle
$$
\mathcal{F} \to Rb_*b^{-1}\mathcal{F} \to Q \to \mathcal{F}[1]
$$
Using that $\mathcal{F} \to b_*b^{-1}\mathcal{F}$ is injective
as before and using what we just said, we see that $Q$ has nonzero
cohomology sheaves only in degrees $0, 1, 2$ sitting on $Z$.
Moreover, these cohomology sheaves are torsion by
Lemma \ref{lemma-torsion-direct-image}.
By induction and the spectral sequence of
Derived Categories, Lemma \ref{derived-lemma-two-ss-complex-functor}
we see that $H^i(X, Q)$ is zero for
$i > 2 + 2\dim(Z) \leq 2 + 2(n - 2) = 2n - 2$. Thus it suffices
to prove that $H^i(X', b^{-1}\mathcal{F}) = 0$ for
$i > 2n$. At this point we use the morphism
$$
f : X' \to \mathbf{P}^1_k
$$
whose fibres have dimension $< n$. Hence by induction we see that
$R^if_*b^{-1}\mathcal{F} = 0$ for $i > 2(n - 1)$.
We conclude by the Leray spectral sequence
$$
H^i(\mathbf{P}^1_k, R^jf_*b^{-1}\mathcal{F})
\Rightarrow
H^{i + j}(X', b^{-1}\mathcal{F})
$$
and the fact that $\dim(\mathbf{P}^1_k) = 1$.
\end{proof}